% Generated draft: requires mathematical and source review.
\documentclass[11pt]{article}
\usepackage[margin=25mm]{geometry}
\usepackage{fontspec}
\setmainfont{Latin Modern Roman}
\usepackage{amsmath,amssymb,mathtools,unicode-math}
\setmathfont{Latin Modern Math}
\usepackage{enumitem,needspace,xurl,hyperref}
\setlength{\parindent}{0pt}
\setlength{\parskip}{5pt}
\setlength{\emergencystretch}{3em}
\newcommand{\N}{\mathbb{N}}
\newcommand{\Z}{\mathbb{Z}}
\newcommand{\Q}{\mathbb{Q}}
\newcommand{\R}{\mathbb{R}}
\newcommand{\C}{\mathbb{C}}
\newcommand{\F}{\mathbb{F}}
\newcommand{\A}{\mathbb{A}}
\newcommand{\PP}{\mathbb{P}}
\newcommand{\E}{\mathbb{E}}
\newcommand{\eps}{\varepsilon}
\newcommand{\dd}{\,\mathrm{d}}
\begin{document}
\section*{UnsolvedMath ID: 4700017}
\textbf{Source problem code:} catalogue entry
\subsection{Definitions and mathematical statement}
Compare planar quadratic polynomial systems $\dot x=P(x,y)$, $\dot y=Q(x,y)$, with $\max(\deg P,\deg Q)=2$, and planar piecewise-linear systems whose two affine-linear vector fields are separated by a straight line. A limit cycle is an isolated periodic orbit. Call it algebraic if it is contained in the zero set of a nonzero real polynomial. A limit cycle is nonhyperbolic if its transverse return multiplier equals one. For a center with a period annulus, a critical period is a periodic orbit at which the period function has zero derivative with respect to a regular transverse orbit parameter; the source does not specify whether multiplicities are counted.

Determine the unknown entries in the source's comparison, including the status and sharpness of the following proposed values, listed as quadratic versus piecewise-linear:
\begin{enumerate}
\item Maximum number of limit cycles: $4?$ versus $3?$.
\item Maximum number of algebraic limit cycles: $1?$ versus $2?$.
\item Maximum number of nonhyperbolic algebraic limit cycles: $0?$ versus $1?$.
\item Coexistence of an algebraic and a nonalgebraic limit cycle in one system: unknown for quadratic systems; described as impossible for the piecewise-linear class.
\item Maximum number of critical periods: $2?$ versus unknown.
\end{enumerate}
Every question mark is retained as a qualification, not converted into a proved upper bound or an attained maximum. The intended extremal class for the piecewise-linear column needs clarification: the supplied summary refers to crossing limit cycles, but the statement itself does not exclude sliding or other nonsmooth behaviors or specify continuity across the line.

The supplied research summary focuses the remaining issues on quadratic coexistence and sharp piecewise-linear critical-period counts. It also calls other entries known and writes quadratic ``$\leq4$'' and piecewise-linear ``$\leq3$'' crossing-cycle bounds. Those assertions are not accompanied by references or hypotheses and are not a replacement for the qualified comparison in the source statement.
\subsection{Short English statement}
Complete the comparison of quadratic and two-zone piecewise-linear planar systems: count limit cycles, algebraic cycles, nonhyperbolic algebraic cycles, and critical periods, and decide coexistence of algebraic and nonalgebraic cycles. Preserve the source's question marks on candidate maxima.
\subsection{Sources}
\begin{itemize}
\item Armengol Gasull, \emph{Some open problems in low dimensional dynamical systems} (2020), problem environment 17, ``Algebraic limit cycles and related questions,'' \url{https://arxiv.org/abs/2012.02524}; accessed 2026-07-29, with source-TeX extraction recorded in the supplied provenance.
\item Local registry \texttt{data/unsolvedmath/problems.json}, permanent UnsolvedMath ID $4700017$, source code catalogue entry, \emph{Algebraic limit cycles and related questions}. Dataset attribution: Ulam AI and the UnsolvedMath Contributors.
\end{itemize}
\end{document}
